\section{Estabilidad al revelar una pieza}
\paragraph{Lema revisado.}
Sean $P$ el punto de torres bistocástico, $\max p\le M_0/k$ y
$\|P-J_k\|_{\op}\le r<1$, con $J_k=\ind\ind^\top/k$. Al quitar una fila y una columna
y hasta dos conjuntos con a lo sumo una casilla por fila y por columna, el nuevo
punto $P'$ en la rama difusa satisface $p'=p(1+O(k^{-1}))$ en el soporte conservado.
Los conjuntos describen las dos diagonales ocupadas por la nueva reina.
Las constantes sólo dependen de $M_0,r$ y $k$ se toma suficientemente grande.
\paragraph{Prueba.}
Sepárese la restricción $q$ de $P$ al soporte nuevo, extendida por cero en sus vetos.
Cada margen conservado pierde a lo sumo $3M_0/k$. En coordenadas de fila y columna
se busca la rama pequeña
\[
 p'_{ij}(1-p'_{ij})=q_{ij}(1-q_{ij})e^{a_i+b_j}.
\]
El mapa de márgenes en $(a,b)=0$ tiene jacobiano de forma
$\sum_{ij}\tau(q_{ij})(a_i+b_j)^2$, donde $\tau(q)=q(1-q)/(1-2q)\ge q$.
En la gauge $\sum a_i=\sum b_j$, prolongar los potenciales por cero en los vértices
eliminados conserva la ortogonalidad al modo nulo. El jacobiano de masas para $P$ tiene
brecha al menos $1-r$. Quitar aristas cuesta $O(k^{-1})\|(a,b)\|_2^2$ pues los pesos
borrados incidentes en cada vértice conservado suman $O(k^{-1})$.
Se obtiene inverso $\ell^2$ acotado. Para un segundo miembro $f$ compatible, resulta
$\|v\|_2\le C\sqrt k\|f\|_\infty$. La ecuación de una coordenada, su diagonal
$1+O(k^{-1})$ y la norma $\ell^2$ de su fila cruzada $O(k^{-1/2})$ implican
$\|v\|_\infty\le C\|f\|_\infty$. La segunda derivada de la rama pequeña por casilla
es $O(q)$ en un entorno fijo. Por sumación en los márgenes, el resto de Newton tiene
constante cuadrática uniforme en supremo. Aplicar el inverso fijo al defecto $O(k^{-1})$
es una contracción en la bola de radio $C'/k$ para $C'$ suficientemente grande y $k$
grande. Como la derivada logarítmica por casilla en el soporte es acotada,
$p'/q=1+O(k^{-1})$. La ecuación asegura la estacionariedad puesto que $q(1-q)$
factoriza sobre el soporte conservado. La concavidad por fila del funcional de torres
identifica el máximo. En efecto, si todos los pesos son $<1/2$ el Hessiano es estrictamente
negativo. Si una entrada $p_0>1/2$, $\sum_{i\ne0}\tau(p_i)\le\tau(1-p_0)=-\tau(p_0)$,
ya que $\tau(t)/t$ es creciente por debajo de $1/2$; Cauchy--Schwarz en la tangente
prueba la concavidad también allí. Los extremos se tratan por continuidad.
\paragraph{Cambio del núcleo centrado.}
Insértese $P'-J_{k-1}$ con ceros en la fila y columna retiradas para compararlo con $P-J_k$.
Las entradas son $O(k^{-1})$, el cambio tiene suma absoluta $O(1)$ y norma de
Hilbert--Schmidt $O(k^{-1/2})$. En efecto, $O(k)$ posiciones se modifican en $O(k^{-1})$
y las restantes en $O(k^{-2})$. Todo el segmento conserva brecha, reducida si es preciso.
